% !PRINTED 139
\weberpage{139}
%<<CONT
1 mit 2 und 4 mit 5 vertauscht, in die erste über, und ebenso lassen sich die drei anderen Annahmen über $b$ auf einander zurückführen. Es bleiben also nur zwei Möglichkeiten zu untersuchen:
\[
\begin{gathered}
a = (1, 2, 3) (4, 5, 6), \quad b = (1, 2, 3, 4, 5) \\
a = (1, 2, 3) (4, 5, 6), \quad b = (1, 2, 3, 5, 4).
\end{gathered}
\]
Davon ist aber die erste zu verwerfen, weil sie in
\[
a^2 b = (1, 4, 6)
\]
einen dreigliedrigen Cyklus ergeben würde, der nicht vorkommen kann. Es bleibt also nur die einzige Möglichkeit:
\[
a = (1, 2, 3) (4, 5, 6), \quad b = (1, 2, 3, 5, 4),
\]
und es giebt also, wenn man äquivalente Gruppen als nicht verschieden betrachtet, nur \gesperrt{eine} einfache Gruppe 60\textsuperscript{sten} Grades.

Diese beiden Elemente $a, b$ kann man als erzeugende Elemente der ganzen Gruppe betrachten und man kann durch ihre wiederholte Zusammensetzung auf sehr mannigfaltige Art die ganze Gruppe vom 60\textsuperscript{sten} Grade herleiten. Wir bekommen ausser dem Einheitselemente die Permutationen 5\textsuperscript{ten} Grades:
\[
\begin{gathered}
(2, 3, 6, 5, 4), \quad (1, 3, 5, 6, 4), \quad (1, 2, 5, 4, 6), \\
(1, 2, 6, 3, 5), \quad (1, 2, 4, 6, 3), \quad (1, 2, 3, 5, 4),
\end{gathered}
\]
die mit ihren Potenzen 24 ergeben; ferner die Permutationen 3\textsuperscript{ten} Grades:
\[
\begin{array}{ll}
(1, 2, 3) (4, 5, 6), & (1, 3, 5) (2, 4, 6), \\
(1, 2, 4) (3, 6, 5), & (1, 3, 6) (2, 4, 5), \\
(1, 2, 5) (3, 6, 4), & (1, 4, 5) (2, 3, 6), \\
(1, 2, 6) (3, 4, 5), & (1, 4, 6) (2, 3, 5), \\
(1, 3, 4) (2, 6, 5), & (1, 5, 6) (2, 3, 4),
\end{array}
\]
die mit ihren Quadraten zusammen 20 ergeben. Endlich erhält man noch 15 Permutationen 2\textsuperscript{ten} Grades:
\[
\begin{array}{lll}
(1, 2) (3, 4), & (1, 2) (5, 6), & (3, 4) (5, 6), \\
(1, 3) (4, 5), & (1, 3) (2, 6), & (2, 6) (4, 5), \\
(1, 4) (2, 5), & (1, 4) (3, 6), & (2, 5) (3, 6), \\
(1, 5) (2, 3), & (1, 5) (4, 6), & (2, 3) (4, 6), \\
(1, 6) (2, 4), & (1, 6) (3, 5), & (2, 4) (3, 5),
\end{array}
\]
von denen je drei in einer Reihe stehenden mit dem Einheitselemente zusammen eine Gruppe 4\textsuperscript{ten} Grades bilden.

Dass diese Gruppe auch dargestellt werden kann durch die Permutationen von fünf Ziffern, ergiebt sich nun schon daraus,
%CONT>>
